\documentclass[12pt]{article}

\usepackage[utf8]{inputenc}
\usepackage[T1]{fontenc}
\usepackage[romanian]{babel}
\usepackage{amsmath, amssymb, amsthm}
\usepackage{geometry}
\geometry{a4paper, margin=2.5cm}

\title{\Large\textbf{Principiul de Reflexie al lui Schwarz}}
\date{}

\theoremstyle{plain}
\newtheorem{proposition}{Propoziție}

\begin{document}

\maketitle

\begin{proposition}[Principiul de reflexie al lui Schwarz]
Fie $\Omega$ o mulțime deschisă din semiplanul superior


\[
\mathbb{H} = \{ z \in \mathbb{C} : \operatorname{Im} z > 0 \},
\]


astfel încât $\Omega$ intersectează axa reală într-un interval deschis $I \subset \mathbb{R}$.

Presupunem că $f : \Omega \cup I \to \mathbb{C}$ este continuă pe $\Omega \cup I$, olomorfă în $\Omega$, iar pe $I$ satisface


\[
f(x) \in \mathbb{R} \quad \text{pentru orice } x \in I.
\]



Definim funcția extinsă


\[
F(z) =
\begin{cases}
f(z), & \operatorname{Im} z \ge 0, \\
\overline{f(\overline{z})}, & \operatorname{Im} z < 0.
\end{cases}
\]



Atunci $F$ este olomorfă în mulțimea simetrică


\[
\Omega \cup I \cup \{ \overline{z} : z \in \Omega \},
\]


și reprezintă reflexia olomorfă a lui $f$ prin axa reală.
\end{proposition}

\end{document}
